% Responsable(s): por asignar
% Objetivo: interpretar los resultados, contrastar ventajas y limitaciones,
% discutir diferencias entre estudios y relacionar los hallazgos con IC-1400.

% Redacción pendiente.
Los resultados analizados muestran que el rendimiento de una memoria no puede
evaluarse utilizando únicamente una característica. Existe un compromiso entre
latencia y ancho de banda, ya que tecnologías como HBM pueden proporcionar un
ancho de banda considerablemente mayor que DDR, pero no necesariamente una
menor latencia \cite{ibeid2025aurora}. Por esta razón, la tecnología más
adecuada depende de las necesidades y del patrón de acceso de cada aplicación.

Además, el rendimiento observado depende tanto de la memoria como de la
arquitectura del sistema. Esmaili-Dokht et al. muestran que sistemas que
utilizan una misma tecnología de memoria pueden presentar diferencias
importantes de latencia debido a factores como la jerarquía de caché, la
interconexión y la microarquitectura \cite{esmailidokht2024mess}. De manera
similar, Velten et al. muestran que la organización de la jerarquía de memoria,
la ubicación de los datos y las características de acceso local o remoto pueden
producir diferencias importantes en latencia y ancho de banda en sistemas de
servidores \cite{velten2022memory}. Las restricciones internas y el controlador
de memoria también pueden limitar el aprovechamiento del ancho de banda
disponible \cite{eyerman2022dram}.

En conjunto, los artículos revisados indican que no existe una tecnología de
memoria que presente el mejor rendimiento en todas las situaciones. Algunas
aplicaciones pueden beneficiarse principalmente de un mayor ancho de banda,
mientras que otras dependen más de reducir los tiempos de acceso. Por ello, la
evaluación de una tecnología de memoria debe considerar conjuntamente la
latencia, el ancho de banda y las características de la arquitectura donde será
utilizada.

Al comparar estas tecnologías, también se observa que la evolución
de DDR4 y DDR5 está orientada principalmente a aumentar la capacidad de 
transferencia. DDR5 puede ofrecer mayor ancho de banda gracias a cambios 
en su organización y a velocidades de transferencia más elevadas 
\cite{steiner2022dramsys,choi2024ddr5}. Sin embargo, esto no significa
que DDR4 deje de ser competitiva, ya que determinadas cargas sensibles a 
la latencia pueden obtener resultados favorables con esta tecnología 
\cite{steiner2022dramsys}. Trabajos como el de Shahbazi y Tsui muestran que 
una mejor organización del controlador puede mejorar notablemente el 
rendimiento de DDR4, incluso sin cambiar la memoria en sí \cite{shahbazi2026ddr4}. 
Esto confirma que el rendimiento final depende tanto de la generación de memoria 
como de su integración y gestión en el sistema.